\section{Markov geometry and the finite-state law}\label{sec:v6-model}
We first specify the geometry independently of the observation apparatus.
Let $A$ be a primitive zero--one matrix on an alphabet $\mathcal A$ of
cardinality $b\ge2$, with Perron root $\lambda_A>1$. Write $\Sigma_A$
for its one-sided shift and $\mathcal L_n$ for its words of length $n$.
The empty word is denoted by $\varnothing$. Forbidden transitions are
permitted; in particular, a vertex may have only one successor.

Let $D\subset\R^d$ be compact, of diameter at most one, and let
$h_i:D\to D$ be injective contractions. Assume that the compact sets
$h_i(D)$ have pairwise distance at least $g>0$. For $v=i_0\cdots i_{n-1}$
write $h_v=h_{i_0}\circ\cdots\circ h_{i_{n-1}}$. We require scales
$r_\varnothing=1$, $0<r_-\le r_i\le r_+<1$, and $D_0\ge1$ such that
\begin{align}
 D_0^{-1}r_v|x-y|&\le |h_vx-h_vy|\le r_v|x-y|,
                                                        \label{eq:v6-metric}\\
 D_0^{-1}r_ur_v&\le r_{uv}\le r_ur_v.                    \label{eq:v6-scales}
\end{align}
These conditions are imposed for words for which the compositions are
used. They hold for similarities in every dimension, and for increasing
one-dimensional $C^{1+\gamma}$ contractions with uniformly positive
derivative and bounded logarithmic distortion, with $r_v=\sup|h_v'|$.
For higher-dimensional conformal maps, \eqref{eq:v6-metric} is an explicit
metric bounded-distortion hypothesis; no global co-Lipschitz conclusion
is inferred merely from a pointwise derivative estimate.

The coding map $\pi:\Sigma_A\to K\subset D$ is injective. Set
$T\pi(\omega)=\pi(\sigma\omega)$ and $K_v=\pi([v])$.
Let $\mu$ be the image of a stationary, fully supported H\"older Gibbs
measure on $\Sigma_A$. Normalize its potential $\varphi$ by
$P_A(\varphi)=0$. In particular, for a fixed $C_G$,
\begin{equation}
 C_G^{-1}e^{S_n\varphi(\omega)}\le\mu[v]\le
 C_Ge^{S_n\varphi(\omega)}\quad(\omega\in[v]).             \label{eq:v6-gibbs}
\end{equation}
We use the classical Gibbs and variational pressure theorems for a
primitive finite shift in the form recorded in \cite{Bowen2008}.
The profile theorem below needs only stationarity and a uniform positive
lower bound on each admissible one-letter conditional probability.

\subsection{Observables and acquisition times}
Fix $0<\beta\le1$, put $\tau=2\beta$, and let $f:K\to\R^m$ satisfy
\begin{align}
 |f(x)-f(y)|&\le H|x-y|^\beta,\label{eq:v6-holder}\\
 \sum_{j=0}^J|f(T^jx)-f(T^jy)|^2&\ge c_o|x-y|^{2\beta}
                \quad(x,y\in K)                         \label{eq:v6-delay}
\end{align}
for some finite $J$ and $c_o>0$. Thus the present output need not identify
the state; finitely many successive outputs must do so quantitatively.
Condition \eqref{eq:v6-delay} concerns the actual observable on $K$, not
an assigned distance on symbolic words.

At a renewal an independent state $U$ of law $\mu$ is observed exactly.
Between renewals the state evolves by $T$. Cycle lengths are independent
copies of a positive integer $L$, independent of the acquired states,
with finite mean. The age weights and their tails are
\begin{equation}
 w_k=\frac{\Pp(L>k)}{\E L},\qquad W(n)=\sum_{k\ge n}w_k,
 \qquad
 0<\underline q\le\frac{w_{k+1}}{w_k}\le\overline q<1.
                                                        \label{eq:v6-weights}
\end{equation}
An initial acquisition occurs at zero; the next renewal time is not
announced. Nongeometric lengths satisfying \eqref{eq:v6-weights} are
allowed. Write $X_t$ for this process.

An observer keeps $I_t$ in a set of at most $M$ elements. Its update may
use $I_{t-1}$, absolute time, the renewal indicator, and the acquired
state at a renewal. The post-update decoder uses $I_t$, absolute time,
and independent coding randomness, but does not separately receive the
renewal indicator or a retained real-valued observation. All
observationally dependent persistent variables are counted in $I_t$.
Stochastic kernels, private randomness, arbitrary fixed real constants,
and transient computation are allowed. Section~\ref{sec:v6-operational}
compares this convention with other finite-memory interfaces.

Set
\begin{align}
 Z_w^f(x)&=(\sqrt{w_k}f(T^kx))_{k\ge0}\in\ell^2(\N_0;\R^m),\nonumber\\
 e_M^2(Z_w^f)&=\inf_{c_1,\ldots,c_M}
       \int\min_{i\le M}\|Z_w^f(x)-c_i\|^2\,d\mu(x),\nonumber\\
 \mathcal R_M&=\inf_{\mathcal C}\limsup_{N\to\infty}\frac1N
       \sum_{t<N}\E|f(X_t)-\widehat f_t|^2.                \label{eq:v6-risk}
\end{align}
Centres in the quantization problem may lie anywhere in the Hilbert
space. The infimum in the last line ranges over the entire preceding
observer class, not just stationary or cylinder observers.

\subsection{A computable geometric profile}
For $v\in\mathcal L_n$ and its suffix $v_{k:}=i_k\cdots i_{n-1}$ define
\begin{equation}
 V_w(v)=\sum_{k<n}w_kr_{v_{k:}}^\tau+W(n),\qquad
 a_w(v)=\mu[v]V_w(v),\qquad a_w(\varnothing)=1.            \label{eq:v6-energy}
\end{equation}
Starting from the root, split a leaf of maximal $a_w$ into all its
admissible one-symbol successors, with a fixed rule for ties. Let
$\mathcal T_N$ be the last tree in this sequence having at most $N$
leaves, let $\mathcal P_N$ be its leaves, and put
\begin{equation}
 G_N(w)=\sum_{v\in\mathcal P_N}a_w(v).                   \label{eq:v6-profile}
\end{equation}
For the finitely many budgets preceding the first root split, the tree
is the root. Unary splits are genuine operations. They do not increase
the number of leaves, but they cannot occur in an infinite uninterrupted
chain. Consequently the last tree in this definition exists.

The construction is a finite mathematical algorithm on cylinder masses,
scales and renewal tails. Effective numerical evaluation additionally
requires certified access to those real inputs; no computability
assertion for arbitrary Borel or H\"older data is implicit.

\begin{theorem}[Markov renewal profile]\label{thm:v6-profile}
Under \eqref{eq:v6-metric}--\eqref{eq:v6-weights}, there are positive
constants $c,C$, independent of $M$, such that
\begin{equation}
 cG_M(w)\le e_M^2(Z_w^f)\le\mathcal R_M\le CG_M(w)
                  \quad(M\ge1).                        \label{eq:v6-profile-law}
\end{equation}
Every randomized time-dependent observer satisfies the middle lower
bound with its lower limiting average loss. A stationary deterministic
observer attains the upper order.

More precisely, the complete admissible tree $\mathcal T_N$ is suffix
closed. If it has $L$ leaves, all its vertices can be retained as states,
and their number is at most
\begin{equation}
                 (b+1)(2L-1).                           \label{eq:v6-state-count}
\end{equation}
Its exact number is $1+\sum_{v\text{ internal}}d^+(v)$, where $d^+(v)$
is the number of children actually present. No time-depth factor is
required. The constants in \eqref{eq:v6-profile-law} may depend on the
graph, distortion, separation, conditional masses, observable, and the
two ratios in \eqref{eq:v6-weights}.
\end{theorem}

For the pressure statement assume additionally that a H\"older potential
$\psi$ satisfies $\log r_v=S_n\psi(\omega)+O(1)$ uniformly on $[v]$,
and that
\begin{equation}
             \lim_{n\to\infty}n^{-1}\log w_n=\log q<0.   \label{eq:v6-tail-rate}
\end{equation}
The inverse-derivative potential gives $\psi$ in the one-dimensional
smooth case. For similarities, $\psi(\omega)=\log r_{\omega_0}$.

\begin{theorem}[Intrinsic pressure exponent]\label{thm:v6-pressure}
There are unique roots $s_g,s_t\in(0,1)$ of
\begin{equation}
 P_A(s_g\varphi+\tau s_g\psi)=0,\qquad
 P_A(s_t\varphi)+s_t\log q=0.                             \label{eq:v6-roots}
\end{equation}
Writing $s_* =\max(s_g,s_t)$, one has
\begin{equation}
 \lim_{M\to\infty}\frac{-\log\mathcal R_M}{\log M}
 =\lim_{M\to\infty}\frac{-\log e_M^2(Z_w^f)}{\log M}
 =\frac{1-s_*}{s_*}.                                    \label{eq:v6-exponent}
\end{equation}
The spatial and survival roots exchange dominance at
$q_c=\exp\{-P_A(s_g\varphi)/s_g\}$. This formula uses the exponential
rate of the renewal tail. Polynomial factors in that tail can change
the critical correction without changing either root.
\end{theorem}
